% !PRINTED 689
\weberpage{689}
%<<CONT
and if $y_k, y_k'$ belong to an imaginary pair,
\weq{9}{z_k = \log y_k y_k',}
by which the definition becomes completely unambiguous.

Now we determine the variables $\xi_1, \xi_2, \ldots, \xi_{\nu-1}, \xi_\nu$ from the linear equations
\weq{10}{\begin{array}{l}
\xi_1 l_{1,1} + \cdots + \xi_{\nu-1} l_{\nu-1,1} + \delta_1 \xi_\nu = z_1\\
\xi_1 l_{1,2} + \cdots + \xi_{\nu-1} l_{\nu-1,2} + \delta_2 \xi_\nu = z_2\\
. \; . \; . \; . \; . \; . \; . \; . \; . \; . \; . \; . \; . \; .\\
\xi_1 l_{1,\nu} + \cdots + \xi_{\nu-1} l_{\nu-1,\nu} + \delta_\nu \xi_\nu = z_\nu,
\end{array}}
whose determinant does not vanish.

If one gives the variables $x_i$ integer values, then by \S.\,188, (2), (3) the $\xi_1, \xi_2, \ldots, \xi_{\nu-1}$ pass over into the system of exponents of the number $\alpha$, and by addition of the equations (10) [\S.\,185, (3)] there results:
\weq{11}{\xi_\nu = \frac{1}{n} \log N_a(\alpha).}

We now pass from the point $(x)$ to the point $(x')$ by setting $x_i' = t^{\frac{1}{n}} x_i$. Thereby $y_i$ passes into
\[ y_i' = t^{\frac{1}{n}} y_i \]
and $z_i$ into
\[ z_i' = z_i + \frac{1}{n} \delta_i \log t. \]

If, then, we put $z_i'$ in place of $z_i$, and denote the values of $\xi_1, \xi_2, \ldots, \xi_\nu$ that then result from (10) by $\xi_1', \xi_2', \ldots, \xi_\nu'$, we have
\weq{12}{\xi_1' = \xi_1, \xi_2' = \xi_2, \ldots, \xi_{\nu-1}' = \xi_{\nu-1}, \xi_\nu' = \xi_\nu + \frac{1}{n} \log t.}

We now propose to delimit the region $S'$ by requiring that
\weq{13}{\begin{array}{c}
0 \leqq \xi_1' < 1, \ldots, 0 \leqq \xi_{\nu-1}' < 1, \\
\xi_\nu' < \frac{1}{n} \log t
\end{array}}
shall hold. By this we achieve that the lattice points $(x')$ lying in the region $S'$ represent only reduced integers $\alpha$, and all and only those whose absolute norm is less than $t$.

% German ends 'die Anzahl der in $\Omega$ enthaltenen |' with the noun on printed 690;
% English word order forces the noun onto this page so that 'contained in $\Omega$' can
% follow it. i0713 opens correspondingly at 'and if we therefore'.
But among the reduced numbers there are always $w$ at a time associated with one another, if $w$ denotes the number of the roots of unity contained in $\Omega$,
%CONT>>
